\chapter{Conclusions}
\label{cap:conclu}

This thesis has systematically investigated the properties of discrete photonic lattices with non-conventional transverse modes, focusing on configurations that include \( P \) modes, \( S\hspace*{-0.3ex}P \)inter-orbital couplings, and non-symmetric evanescent couplings. The research combined analytical tools, numerical simulations, and experimental validation to address different manifestations of the orbital structure of guided light.

In Chapter~\ref{cap:invisibility}, a honeycomb-type photonic lattice composed of \( p_y \) modes was characterized, demonstrating that the \textit{invisibility angle} induces a quasi-flat band regime around \( \theta \approx 0.56 \) radians. Through a detailed analysis incorporating non-orthogonality corrections and long-range couplings, it was established that Aharonov-Bohm-type confinement persists even in the presence of these perturbations. In particular, it was confirmed that non-orthogonality effects are essential to correctly describe localization at the lattice edges, which was quantified using the inverse participation ratio (IPR) and experimentally corroborated.

Chapter~\ref{cap:asymmetric} addressed the phenomenon of non-symmetric evanescent coupling in photonic dimers with different refractive index contrasts. A generalization of coupled-mode theory was presented, including non-orthogonality terms, which allows correct modeling of the observed directional dynamics. Experimentally, an energy transfer imbalance dependent on the input channel was verified, in agreement with theoretical predictions based on non-symmetric effective coupling matrices.

In Chapter~\ref{cap:molecules}, the implementation of SP-SSH lattices constructed from photonic molecules was studied. It was demonstrated that structural dimerization, combined with hybridization between \( S \) and \( P \) modes, gives rise to two independently controllable topological phase transitions. Phases with different quantized bulk polarizations were identified, and an experimental protocol was designed to detect edge modes through the ratio \( I_{\text{borde}} / I_{\text{total}} \)in excellent agreement with numerical simulations. This model extends the physics of the one-dimensional SSH system to the multi-orbital context, revealing new mechanisms of topological localization mediated by internal degrees of freedom.

From a methodological perspective, this thesis developed a continuous numerical method based on eigenmode expansion (EME), which extends coupled-mode theory to the regime of closely spaced waveguides or arbitrary refractive index contrasts. This approach was complemented by experimental techniques for fabricating and characterizing multi-orbital lattices, enabling the calibration of propagation constants for distinct guided modes.

The obtained results open new possibilities for designing photonic devices that exploit localization and control mechanisms based on geometry, transverse symmetry, and topology. Furthermore, a theoretical-experimental framework adaptable to more complex platforms is established, for example, in systems with orbital angular momentum \citep{OAMCaging}.
